\documentclass[10pt,a4paper]{amsart}
\usepackage[margin=19mm]{geometry}
\usepackage{amsmath,amssymb,mathtools}
\usepackage[scr=rsfs,cal=euler]{mathalfa}
\usepackage{xfrac}
\usepackage[unicode,hidelinks,bookmarks=false]{hyperref}
\newcommand{\ie}{i.e.\ }
\newcommand{\cf}{cf.\ }
\newcommand{\resp}{respectively\ }
\newcommand{\Subsection}{\subsection}
\providecommand{\nobreakdash}{\nobreak\hbox{-}}
\setlength{\emergencystretch}{2em}
\multlinegap0pt
\usepackage{bm}
\usepackage{bbm}
\usepackage{dsfont}
\usepackage{xspace}
\usepackage{cancel}
\usepackage{mathtools}
\usepackage{amsthm}
\makeatletter
\@ifundefined{theorem}{\newtheorem{theorem}{Theorem}}{}
\makeatother
\def \codim {\operatorname{codim}}
\def \det {\operatorname{det}}
\def \p {\mathbb{P}}
\def \ver {\operatorname{Vert}}
\title[On quadratic persistence and Pythagoras numbers of totally r]{On quadratic persistence and Pythagoras numbers of totally real projective varieties}
\author{Jong In Han, Jaewoo Jung, Euisung Park}
\date{Source: arXiv:2506.13247v1 (CC BY 4.0)}
\begin{document}
\maketitle

\noindent\emph{Source and licence.} On quadratic persistence and Pythagoras numbers of totally real projective varieties. Version arXiv:2506.13247v1 (CC BY 4.0), distributed under the Creative Commons Attribution 4.0 International licence, as recorded in the arXiv licence metadata for this exact version and shown on the abstract page https://arxiv.org/abs/2506.13247v1. Full rights notice: licenses/arxiv-2506.13247-CC-BY-4.0.txt.\par\smallskip

\noindent\textit{Editorial scope.} Counted unit: the Theorem whose source label is \texttt{lp\_2c+3} in the article (source0.tex lines 504--508, source label \texttt{lp\_2c+3}), delivered together with the article's own complete original proof of it (source0.tex lines 509--553, one \texttt{proof} environment of 45 lines). No mathematics is written, repaired or re-derived: statement and proof are the article's own text line for line. Required context retained uncounted: none. Every definition, display and label of those lines is retained unchanged. Omitted from this package and not redistributed: 1--503, 554--735. Nothing inside a retained excerpt is omitted or altered: every retained body is delivered exactly as the article's lines are written, so no omission receipt of any kind accompanies this package. Citations: the retained excerpts carry no citation command at all, so no bibliography entry of the article is transcribed and no reference is redistributed. Reference adaptation: none was needed and none was made; every reference inside the delivered unit resolves inside it, so no unresolved reference or citation is produced. Printed slips kept literal: no printed word, symbol, hypothesis or number of the counted statement or of its proof was altered. Layout and class: the article's own class is \texttt{amsart} with packages including amsmath, amsthm, amssymb, tikz-cd, comment, hyperref, cleveref, breqn, float, caption, subfig, ytableau, tikz, cases; the wrapper is the lane's pinned \texttt{amsart} editorial wrapper at 10pt on A4, which restates the theorem-like environments the retained text needs and adds the article's own macro definitions that the retained text needs (\texttt{\textbackslash codim}, \texttt{\textbackslash det}, \texttt{\textbackslash p}, \texttt{\textbackslash ver}), with extra wrapper packages (bm, bbm, dsfont, xspace, cancel, mathtools). The delivered numbering is the unit's own. Assets: no retained span contains a figure environment, a graphics-inclusion command, a PDF-page inclusion, a table or a TikZ picture; no image file is copied into this package, the package declares no asset dependency and no image file is redistributed with it. Licence: version arXiv:2506.13247v1 of this article is distributed under the Creative Commons Attribution 4.0 International licence, as recorded in the arXiv licence metadata for this exact version and shown on the abstract page https://arxiv.org/abs/2506.13247v1. A full rights notice is delivered at licenses/arxiv-2506.13247-CC-BY-4.0.txt. Characters: every retained excerpt is pure ASCII, so the delivered engine, XeLaTeX with its default Latin Modern text font, reports no missing character.
\begin{theorem}\label{lp_2c+3}
    Let $X\subset\p^r$ be a nondegenerate projective variety of dimension $n$, codimension $c$, and degree $d\geq 2c+3$. 
    Let $X_q\subset\p^{r-1}$ denote the inner projection of $X$ from a general point $q\in X$.
    If $X_q$ is contained in a variety of minimal degree as a divisor, then so is $X$.
\end{theorem}

\begin{proof}
    Denote $p_0:=q$ and $Y_1:=Y$.
    Let $p_1$ be a general point of $Y_1$.
    Denote the projection of $Y_1\subset\p^{r-1}$ from $p_1$ by $Y_2\subset\p^{r-2}$, and let $p_2$ be a general point of $Y_2$.
    By repeating this, we get $p_0,p_1,\cdots,p_{c-2}$ and $Y_{c-1}=\p^{n+1}$.
    We pick $n+2$ general points in $Y_{c-1}=\p^{n+1}$ and denote them by $p_{c-1},p_c,\cdots,p_r$.
    Then $p_0,p_1,\cdots,p_r$ are linearly independent.
    Let $\{x_i = p_i^* \mid 0 \leq i \leq r \}$ be the dual basis of $\{p_0,p_1,\cdots,p_r\}$.
    Thus $x_0,x_1,\cdots, x_r$ form a homogeneous coordinate system of $\p^r$.
    Denote by $S$ and $R$ the homogeneous coordinate rings of $\p^r$ and $\p^{r-1}$, respectively.
    Let $J\subset S$ be the ideal generated by $I(X)_2$ and $W\subset\p^r$ be the scheme defined by $J$.
    Then $$(J\cap R)_2=(I(X)\cap R)_2=I(X_q)_2.$$
    As $d\geq 2c+3$, we have $\deg X_q\geq 2\codim X_q+4$. 
    Hence $\dim I(X_q)_2\leq \binom{\codim X_q}{2}$. 
    However, as $X_q$ is contained in a variety of minimal degree as a divisor, it holds that 
    $$\dim I(X_q)_2=\binom{\codim X_q}{2}\text{, and therefore}~ I(X_q)_2=I(Y)_2.$$
    In summary, we have $(J\cap R)_2=I(Y)_2$ and $J\cap R\supset I(Y)$.
    
    As $p_0$ is not contained in the vertex set $\ver(W_{\text{red}})$ of $W_{\text{red}}$, there exists $f_0\in J_2$ such that $\frac{\partial f_0}{\partial x_0}\neq 0$.
    Also as $p_1\notin\ver(Y)\subset\p^{r-1}$, there exists $f_1\in I(Y)_2$ such that $\frac{\partial f_1}{\partial x_1}\neq 0$. Since $f_1\in R$, we have $\frac{\partial f_1}{\partial x_0}=0$.
    Let $Y_{p_1}\subset\p^{r-2}$ be the inner projection of $Y$ from $p_1$.
    As $p_2\notin\ver(Y_{p_1})$, there exists $f_2\in I(Y_{p_1})_2$ such that $\frac{\partial f_2}{\partial x_2}\neq 0$ while $\frac{\partial f_2}{\partial x_0}=\frac{\partial f_2}{\partial x_1}=0$.
    By repeating this, we get $f_0,f_1,\cdots,f_{c-2}\in J_2$ such that $\frac{\partial f_i}{\partial x_i}\neq 0$ while $\frac{\partial f_i}{\partial x_j}=0$ for all $j<i$.
    Hence by taking a proper generating set of $J$, one can regard the Jacobian matrix of $W\subset\p^r$ contains the following form of $(c-1) \times (c-1)$ submatrix where $*$ represents a nonzero linear form.
    \[
    M=
    \begin{bmatrix}
        * & 0 & 0 & 0 & 0 \\
        ? & * & 0 & 0 & 0 \\
        ? & ? & * & 0 & 0 \\
        ? & ? & ? & * & 0 \\
        ? & ? & ? & ? & * \\
    \end{bmatrix}
    \]
    Note that $\dim W = n$ or $n+1$. Define the locus $U:=\{x\in W_{\text{red}}\mid \dim T_xW\geq n+2\}$.
    Then $U\subset V(\det M)$ as sets.
    As $\det M$ is decomposable to linear forms, every irreducible component of $U$ is degenerate.
    Note that the scheme $W$ and the locus $U$ does not depend on the choice of $q\in X$.
    Therefore, $q\notin U$ and $\dim T_qW\leq n+1$.
    Hence $$\dim (K_1(I(X),q))_1=\dim (K_1(J,q))_1=\codim T_qW\geq c-1.$$
    Since $\dim I(X)_2=\dim I(X_q)_2+\dim (K_1(I(X),q))_1$, we obtain
    $$\dim I(X)_2\geq\binom{e-1}{2}+c-1=\binom{c}{2}.$$
    Since we have $\dim I(X)_2\leq \binom{c}{2}$ from the condition, $$\dim I(X)_2=\binom{c}{2}.$$
    Together with $d\geq 2c+3$, we conclude that $X$ is contained in a variety of minimal degree as a divisor.
\end{proof}

\end{document}
